\documentclass[tikz,border=0pt]{standalone}
% Shared visual language for every block diagram in figures/diagrams.
% Build with XeLaTeX (see figures/diagrams/build.sh).
\usepackage{fontspec}
\setmainfont{Inter}[Numbers=Lining]
\setmonofont{DejaVu Sans Mono}[Scale=0.86]
\usepackage{xcolor}
\usetikzlibrary{arrows.meta,positioning,calc,fit,backgrounds,shapes.geometric,
  shapes.misc,matrix,decorations.pathreplacing,patterns,quotes}

% Palette (matches python/figstyle.py) --------------------------------------
\definecolor{surface}{HTML}{FCFCFB}
\definecolor{ink}{HTML}{0B0B0B}
\definecolor{ink2}{HTML}{52514E}
\definecolor{muted}{HTML}{898781}
\definecolor{hair}{HTML}{E1E0D9}
\definecolor{memline}{HTML}{2A78D6}  \definecolor{memfill}{HTML}{E7F0FC}
\definecolor{aluline}{HTML}{14875E}  \definecolor{alufill}{HTML}{E3F5EE}
\definecolor{ctlline}{HTML}{B07800}  \definecolor{ctlfill}{HTML}{FDF2D9}
\definecolor{regline}{HTML}{6250D6}  \definecolor{regfill}{HTML}{ECEAFB}
\definecolor{wire}{HTML}{3D3C39}
\definecolor{accent}{HTML}{EB6834}   \definecolor{accentfill}{HTML}{FDE8DF}
\definecolor{pos}{HTML}{E34948}      \definecolor{posfill}{HTML}{FBE1E0}
\definecolor{neg}{HTML}{2A78D6}      \definecolor{negfill}{HTML}{DCE9FA}
\definecolor{zerofill}{HTML}{F0EFEC}

\tikzset{
  every picture/.style={line cap=round, line join=round, font=\small, text=ink},
  % canvas
  canvas/.style={show background rectangle,
    background rectangle/.style={fill=surface, draw=hair, rounded corners=8pt, line width=0.6pt},
    inner frame xsep=16pt, inner frame ysep=14pt},
  ttl/.style={font=\bfseries\large, anchor=north west, inner sep=0pt, text=ink},
  sub/.style={font=\small, anchor=north west, inner sep=0pt, text=ink2},
  % blocks
  blk/.style={rounded corners=3pt, line width=0.8pt, align=center, minimum height=9mm,
    inner xsep=5pt, inner ysep=3pt},
  mem/.style={blk, draw=memline, fill=memfill},
  alu/.style={blk, draw=aluline, fill=alufill},
  ctl/.style={blk, draw=ctlline, fill=ctlfill},
  reg/.style={blk, draw=regline, fill=regfill, rounded corners=1pt},
  io/.style={blk, draw=muted, fill=white},
  ghost/.style={blk, draw=muted, dashed, fill=surface},
  % arithmetic primitives
  add/.style={circle, draw=aluline, fill=white, line width=0.8pt, minimum size=6.2mm,
    inner sep=0pt, font=\normalsize},
  shf/.style={rectangle, rounded corners=1.5pt, draw=aluline, fill=white, line width=0.7pt,
    font=\ttfamily\footnotesize, inner xsep=2.5pt, inner ysep=2pt, minimum height=5mm},
  gain/.style={isosceles triangle, isosceles triangle apex angle=60, draw=aluline, fill=white,
    line width=0.7pt, inner sep=1.2pt, font=\footnotesize},
  dn/.style={circle, draw=wire, fill=white, line width=0.7pt, minimum size=8mm, inner sep=0pt,
    font=\footnotesize},
  mux/.style={trapezium, trapezium angle=72, shape border rotate=270, draw=ctlline, fill=ctlfill,
    line width=0.8pt, minimum height=13mm, minimum width=5mm, inner sep=1pt, font=\scriptsize},
  net/.style={draw=memline, fill=memfill, rounded corners=1.5pt, font=\footnotesize, inner xsep=3pt,
    inner ysep=1.6pt, minimum width=5.5mm, minimum height=4.8mm},
  % wires
  sig/.style={-{Latex[length=2.2mm,width=1.6mm]}, draw=wire, line width=0.75pt},
  bus/.style={-{Latex[length=2.6mm,width=2mm]}, draw=wire, line width=1.5pt},
  ctrl/.style={-{Latex[length=1.8mm,width=1.4mm]}, draw=ctlline, line width=0.65pt, dash pattern=on 3pt off 2pt},
  nohead/.style={-, draw=wire, line width=0.75pt},
  dot/.style={circle, fill=wire, inner sep=0pt, minimum size=3.2pt},
  % labels
  lbl/.style={font=\footnotesize, text=ink2, fill=surface, inner sep=1.2pt},
  lblc/.style={font=\footnotesize, text=ink2, inner sep=1.2pt},
  note/.style={font=\footnotesize, text=ink2, align=left, inner sep=0pt},
  head/.style={font=\bfseries\small, text=ink, inner sep=0pt},
  code/.style={font=\ttfamily\footnotesize, text=ink},
  group/.style={draw=hair, fill=white, rounded corners=6pt, line width=0.8pt, inner sep=7pt},
  grouplabel/.style={font=\bfseries\footnotesize, text=muted, anchor=south west, inner sep=2pt},
}

% Layers: groups sit above the canvas background and below the main layer.
\pgfdeclarelayer{groups}
\pgfsetlayers{background,groups,main}

% A register: box with a clock wedge at the bottom edge.
\newcommand{\clk}[1]{\draw[regline, line width=0.6pt] ($(#1.south)+(-1.4mm,0)$) -- ($(#1.south)+(0,1.6mm)$) -- ($(#1.south)+(1.4mm,0)$);}
% Bus width tick: \bw{(point)}{width}
\newcommand{\bw}[2]{\draw[wire, line width=0.6pt] ($#1+(-0.9mm,-1.2mm)$) -- ($#1+(0.9mm,1.2mm)$); \node[font=\scriptsize, text=ink2, above=0.4mm, inner sep=0.5pt] at #1 {#2};}
\newcommand{\var}[1]{\textit{#1}}
\newcommand{\usub}[1]{\ifcase#1 ₀\or ₁\or ₂\or ₃\or ₄\or ₅\fi}
\AtBeginDocument{\hyphenpenalty=10000\exhyphenpenalty=10000}

\begin{document}
\begin{tikzpicture}[canvas]
  \node[ttl] at (0,0) {Multilevel (Mallat) decomposition};
  \node[sub] at (0,-0.62) {Each level re-decomposes only the LL band. Three levels give ten bands packed in place in one \var{H} × \var{W} array.};
  % chain ---------------------------------------------------------------
  \def\cy{-2.55}
  \node[io, minimum width=16mm] (x) at (0.95,\cy) {\var{x}\\\footnotesize \var{H} × \var{W}};
  \node[alu, minimum width=19mm, minimum height=13mm] (l1) at (3.75,\cy) {2-D DWT\\\footnotesize level 1};
  \node[mem, minimum width=14mm] (ll1) at (6.45,\cy) {LL₁\\\footnotesize \var{H}/2 × \var{W}/2};
  \node[alu, minimum width=19mm, minimum height=13mm] (l2) at (9.15,\cy) {2-D DWT\\\footnotesize level 2};
  \node[mem, minimum width=14mm] (ll2) at (11.85,\cy) {LL₂\\\footnotesize \var{H}/4 × \var{W}/4};
  \node[alu, minimum width=19mm, minimum height=13mm] (l3) at (14.55,\cy) {2-D DWT\\\footnotesize level 3};
  \node[mem, minimum width=14mm] (ll3) at (17.25,\cy) {LL₃\\\footnotesize \var{H}/8 × \var{W}/8};
  \foreach \a/\b in {x/l1, l1/ll1, ll1/l2, l2/ll2, ll2/l3, l3/ll3} \draw[sig] (\a) -- (\b);
  \foreach \l/\j/\s in {l1/1/2, l2/2/4, l3/3/8} {
    \node[mem, minimum width=33mm, minimum height=7.5mm, below=7mm of \l, font=\footnotesize] (d\j) {LH\usub{\j} · HL\usub{\j} · HH\usub{\j}};
    \draw[sig] (\l) -- (d\j);
    \node[note, below=0.8mm of d\j] {each \var{H}/\s{} × \var{W}/\s};
  }
  % packed map ---------------------------------------------------------
  \def\M{6.0} \def\ox{19.6} \def\oy{-1.35}
  \tikzset{band/.style={draw=surface, line width=1.6pt, font=\footnotesize, align=center}}
  \fill[memfill] (\ox,\oy) rectangle ++(\M,-\M);
  \foreach \j/\f in {1/1, 2/0.5, 3/0.25} {
    \pgfmathsetmacro{\h}{\M*\f/2}
    \node[band, fill=negfill, minimum width=\h cm, minimum height=\h cm, anchor=north west, inner sep=0pt] at (\ox+\h,\oy) {LH\usub{\j}};
    \node[band, fill=accentfill, minimum width=\h cm, minimum height=\h cm, anchor=north west, inner sep=0pt] at (\ox,\oy-\h) {HL\usub{\j}};
    \node[band, fill=alufill, minimum width=\h cm, minimum height=\h cm, anchor=north west, inner sep=0pt] at (\ox+\h,\oy-\h) {HH\usub{\j}};
  }
  \node[band, fill=zerofill, minimum width=0.75cm, minimum height=0.75cm, anchor=north west, inner sep=0pt, font=\scriptsize] at (\ox,\oy) {LL₃};
  \draw[hair, line width=0.6pt] (\ox,\oy) rectangle ++(\M,-\M);
  \node[head, anchor=south west] at (\ox,\oy+0.12) {Packed output, row-major};
  \draw[decorate, decoration={brace, amplitude=4pt, mirror}, draw=ink2] (\ox,\oy-\M-0.12) -- node[note, below=5pt] {\var{W}} ++(\M,0);
  % notes ----------------------------------------------------------------
  \node[note, anchor=north west, text width=17.8cm] at (0.2,-6.35) {%
    Work per level falls by 4: the pyramid touches \var{HW}(1 + ¼ + ¹⁄₁₆) samples for three levels, less than 4\var{HW}/3 for any depth.
    The coefficient count stays \var{H}\var{W}, so a \var{J}-level pyramid is an in-place rewrite of one frame buffer.
    The RTL computes one level; deeper pyramids feed LL back through the same engine.};
\end{tikzpicture}
\end{document}
